% !TEX root = ../agr-paper.tex
% ---------------------------------------------------------------------
% ~900 words, compressed from the thesis background chapter.
%
% Four movements: static graph retrieval, agentic navigation, post-hoc
% verification, and the attribution gap. The KBQA-o1 paragraph is load
% bearing twice over -- it is the closest system, and its MCTS ablation
% (GrailQA F1 78.5 -> 48.5) is the field's own evidence that component
% ablation is informative. Citing it as motivation rather than only as a
% competitor is what makes Section 6 look like method rather than
% apology.
%
% RoG gets one extra sentence here because the environment is built
% from its published subgraphs; the direct comparison lives in the
% results section (sec:rog), not here.
%
% CUT from the thesis: the withdrawn "MCTS-inspired" self-correction.
% That is a thesis integrity note. One clause in Section 3 naming the
% search as guided best-first covers what a reader needs. Also cut: the
% thesis's claim that two published papers disagree about ToG's CWQ
% score. Checked in September 2026 against both papers: the 58.9 / 69.5
% pair is ToG-R, a variant from the ToG paper itself, and PoG quotes
% ToG's own 57.1 / 67.6 faithfully. There is no disagreement to cite.
% ---------------------------------------------------------------------

\section{Related Work}\label{sec:related}

Three lines of work bear on question answering over knowledge graphs. We
take them in the order the problem forces: what retrieval fixed before
reasoning can reach, what interleaving retrieval with reasoning buys,
and what checking an answer against evidence adds. A fourth subsection
names something none of them establishes: which component of an agentic
system is responsible for its behaviour. That gap is what this paper
addresses.

\subsection{Static Retrieval over Graphs}\label{sec:rel-static}

Retrieval-augmented generation~\cite{lewis2020retrieval} attaches a
language model to an external corpus. The system embeds the question,
retrieves lexically or semantically similar passages, and lets the model
answer from them. Where the answer sits in a single retrievable passage
this works well. Where it must be composed across passages, the ranking
signal measures the wrong thing. The passage carrying the intermediate
entity resembles neither the question nor anything the system yet knows
to look for. Surveys of the graph--language-model intersection identify
this as the central limitation of flat retrieval~\cite{li2023survey}.

Knowledge graphs make the missing structure explicit. They store facts
as typed edges between identified entities, so composition becomes a
traversal rather than an inference over prose.
GraphRAG~\cite{edge2024local} builds a graph over a corpus and retrieves
structured context in place of passages. Related systems extend the idea
to long-term memory~\cite{gutierrez2024hipporag} and to
temporally-qualified facts~\cite{li2025tgrag}. All of them retrieve
\emph{statically}: a subgraph is fetched around the linked entities in
one shot and handed to the model. Retrieval runs first and runs once. So
it cannot know what the second hop will need, and a fixed-radius
neighbourhood is either too small to contain the answer or so large that
the relevant edges are lost among thousands of irrelevant ones.

\subsection{Agentic Navigation}\label{sec:rel-agentic}

Interleaving retrieval with reasoning removes that constraint. The model
takes one step, observes what it found, and decides where to look next.
Retrieval becomes navigation, and the system becomes an agent operating
on the graph rather than a pipeline reading from it. This inherits the
tool-use and self-reflection patterns established for language agents
generally~\cite{yao2023react,schick2023toolformer,shinn2023reflexion}.

Think-on-Graph~\cite{sun2024think} established the pattern for knowledge
graphs with model-guided beam search over relations and entities.
Plan-on-Graph~\cite{chen2024plan} added two mechanisms to it:
decomposition into sub-objectives, and backtracking driven by
reflection. Reasoning-on-Graphs (RoG)~\cite{luo2024reasoning} takes a
different route. A fine-tuned model generates grounded relation paths as
explicit plans, and the system then traverses them. Because its answers
are read off retrieved paths, they are faithful by construction. RoG is
also the one surveyed system whose data distribution this paper shares:
the knowledge environment of \Cref{sec:environment} is built from its
published subgraphs. \Cref{sec:rog} therefore compares the two directly,
published figures and all. Other systems add tool
abstractions~\cite{jiang2024kgagent}, cope with incomplete
graphs~\cite{xu2024generate}, or generalise across structured
sources~\cite{jiang2023structgpt}. Recent surveys treat the convergence
as a research programme in its own
right~\cite{pan2024unifying,ma2025unifying}.

KBQA-o1~\cite{luo2025kbqao1} is the closest system to AGR's search
strategy, and the most important to distinguish from it. It performs
agentic knowledge base question answering with full Monte Carlo Tree
Search: a policy model proposes actions, a reward model scores them, and
rollouts propagate their values back through the tree. That search wraps
a ReAct-style agent generating logical forms against Freebase, and the
exploration also produces annotations for incremental fine-tuning. AGR
does none of this. Its search is guided best-first with a bounded beam
and backtracking. It carries neither rollouts nor value backpropagation.
The two are separable on that basis, and we name the difference rather
than blur it.

\subsection{Verification and Self-Correction}\label{sec:rel-verify}

A parallel line checks generated content after the fact.
Chain-of-Verification~\cite{dhuliawala2024chain} has a model draft a
response, plan verification questions about it, answer those
independently, and regenerate. Multi-agent debate~\cite{du2024improving}
reaches a related effect by having several instances critique one
another. Both establish the premise this paper's verification layer
rests on: decompose a draft, check its parts, and factual error goes
down. But both check against the model's own knowledge rather than an
external symbolic source, and the limits of models correcting themselves
without such a source are documented~\cite{huang2024selfcorrect}. Hu et
al.~\cite{hu2025selfcorrecting} demonstrate a self-correcting agentic
graph pipeline in a clinical setting. FActScore~\cite{min2023factscore}
decomposes generations into atomic facts for evaluation, which is the
same decomposition applied as a metric rather than as a control step.

None of them provides a claim-level check against the \emph{traversed}
triples, performed \emph{before} the answer is emitted, inside a
graph-navigating agent. That is the gap \Cref{sec:verification} fills.

\subsection{The Attribution Gap}\label{sec:rel-attribution}

A second gap is methodological. Agentic systems are assemblies of
decomposition, scoring, backtracking, and self-correction, and the
literature reports them as wholes. When such a system outperforms a
static baseline, the published work rarely establishes which mechanism
produced the gain, or whether some contribute nothing. Later work then
inherits the entire assembly on faith, along with its cost.

The exception is instructive. KBQA-o1's own ablation removes MCTS while
leaving every other component intact, and GrailQA F1 drops from $78.5$
to $48.5$~\cite{luo2025kbqao1}. That is a thirty-point attribution to a
single mechanism, obtained precisely because the authors ablated rather
than reporting the assembly whole. \Cref{sec:attribution} applies the
same instrument to AGR. It reports the components for which it detects
nothing at the same length as the one for which it does.
